% sec8_5.tex — Section 8.5: FIML

As seen in Chapter~4, the multivariate regression model is a specialization of the
seemingly unrelated regression (SUR) model, which in turn is a specialization of
the multiple-equation system to which three-stage least squares (3SLS) is applicable. The two-stage least squares (2SLS) estimator is the GMM estimator when
each equation of the system is estimated separately, while 3SLS is the GMM estimator when the system as a whole is estimated. This section presents the ML
counterparts of 2SLS and 3SLS.

\subsection*{The Multiple-Equation Model with Common Instruments Restated}

To refresh your memory, the model is an $M$-equation system described by Assumptions~4.1--4.5 and~4.7, with $\mathbf{x}_{tm} = \mathbf{x}_t$ for all $m = 1, 2, \ldots, M$ (so the set of instruments is the same across equations). Still maintaining the convention of denoting the true parameter value by subscript~0, we can write the $M$ equations of the
system as
\begin{equation}
  y_{tm} = \underset{(1 \times L_m)}{\mathbf{z}_{tm}'} \;
  \underset{(L_m \times 1)}{\boldsymbol{\delta}_{0m}}
  + \varepsilon_{tm}
  \quad (m = 1, 2, \ldots, M;\; t = 1, 2, \ldots, n).
  \label{eq:8.5.1} \tag{8.5.1}
\end{equation}
The model has the following features.
\begin{itemize}
\item Unlike in the multivariate regression model, for each equation~$m$, the regressors $\mathbf{z}_{tm}$ may not be orthogonal to the error term, but there are available $K$ predetermined variables $\mathbf{x}_t$ that are known to satisfy the orthogonality conditions
  $\E(\mathbf{x}_t \cdot \varepsilon_{tm}) = \mathbf{0}$ for all~$m$ (this is Assumption~4.3). Defining the $M \times 1$ vector
  $\boldsymbol{\varepsilon}_t$ as
  \begin{equation}
    \underset{(M \times 1)}{\boldsymbol{\varepsilon}_t}
    = \begin{bmatrix} \varepsilon_{t1} \\ \varepsilon_{t2} \\ \vdots \\ \varepsilon_{tM} \end{bmatrix},
    \label{eq:8.5.2} \tag{8.5.2}
  \end{equation}
  the orthogonality conditions can be expressed as
  \begin{equation}
    \E(\mathbf{x}_t \boldsymbol{\varepsilon}_t') = \underset{(K \times M)}{\mathbf{0}}.
    \label{eq:8.5.3} \tag{8.5.3}
  \end{equation}
  These $K$ predetermined variables serve as the common set of instruments in the
  GMM estimation.

\item The \textbf{endogenous variables} of the system are those variables (apart from the
  errors) in the system that are not included in $\mathbf{x}_t$.

\item The rank condition for identification (Assumption~4.4) is that the $K \times L_m$ matrix
  $\E(\mathbf{x}_t\mathbf{z}_{tm}')$ be of full column rank for all~$m$. Equation~$m$ is said to be \textbf{overidentified} if the rank condition is satisfied and $L_m < K$.

\item The error vector is conditionally homoskedastic in that
  $\E(\boldsymbol{\varepsilon}_t\boldsymbol{\varepsilon}_t' \mid \mathbf{x}_t) = \boldsymbol{\Sigma}_0$, and
  $\boldsymbol{\Sigma}_0$ is positive definite (Assumption~4.7).

\item $\E(\mathbf{x}_t\mathbf{x}_t')$ is nonsingular (a consequence of conditional homoskedasticity and the
  nonsingularity requirement in Assumption~4.5).
\end{itemize}

The GMM estimator for this system as a whole is the 3SLS estimator given in
(4.5.12), while the single-equation GMM estimator of each equation in isolation is
2SLS given in (3.8.3).

The following examples will be used to illustrate the concepts to be introduced
shortly.

\begin{example}[A model of market equilibrium]\label{ex:8.1}
An expanded version of
Working's example considered in Section~3.1 is a two-equation system:
\begin{align}
  q_t &= \gamma_{11} p_t + \beta_{11} + \beta_{12} a_t + \varepsilon_{t1}
  \quad \text{(demand)},
  \label{eq:8.5.4} \tag{8.5.4} \\
  q_t &= \gamma_{21} p_t + \beta_{21} + \beta_{22} w_t + \varepsilon_{t2}
  \quad \text{(supply)}.
  \label{eq:8.5.5} \tag{8.5.5}
\end{align}
In this system, $q_t$ is the amount of coffee transacted and $p_t$ is the coffee price.
The variable $a_t$ appearing in the demand equation is a taste shifter, while the
variable $w_t$ is the amount of rainfall that affects the supply of coffee. The
system can be cast in the general format by setting
\[
  y_{t1} = q_t, \; y_{t2} = q_t, \quad
  \mathbf{z}_{t1} = \begin{bmatrix} p_t \\ 1 \\ a_t \end{bmatrix}, \quad
  \mathbf{z}_{t2} = \begin{bmatrix} p_t \\ 1 \\ w_t \end{bmatrix},
\]
\[
  \boldsymbol{\delta}_{01} = \begin{bmatrix} \gamma_{11} \\ \beta_{11} \\ \beta_{12} \end{bmatrix}, \quad
  \boldsymbol{\delta}_{02} = \begin{bmatrix} \gamma_{21} \\ \beta_{21} \\ \beta_{22} \end{bmatrix}.
\]
(Note that $y_{t1}$ and $y_{t2}$ are the same variable.) If $\E(\varepsilon_{tj}) = 0$,
$\E(a_t \varepsilon_{tj}) = 0$, and
$\E(w_t \varepsilon_{tj}) = 0$ for $j = 1, 2$, then $(1, a_t, w_t)$ can be included in the common
set of instruments. So
\[
  \mathbf{x}_t = \begin{bmatrix} 1 \\ a_t \\ w_t \end{bmatrix}.
\]
The other nonerror variables in the system, $p_t$ and $q_t$, are endogenous variables. Provided that the rank condition for identification is satisfied, each
equation is just identified because the number of right-hand side variables
equals that of instruments.
\end{example}

\begin{example}[Wage equation]\label{ex:8.2}
In the empirical exercise of Chapter~3,
we estimated a standard wage equation. Consider supplementing the wage
equation by an equation explaining $KWW$ (the score on the ``Knowledge of
the World of Work'' test):
\begin{align}
  LW_t &= \gamma_{11} S_t + \beta_{11} + \beta_{12} IQ_t + \varepsilon_{t1},
  \label{eq:8.5.6} \tag{8.5.6} \\
  KWW_t &= \gamma_{21} S_t + \beta_{21} + \varepsilon_{t2},
  \label{eq:8.5.7} \tag{8.5.7}
\end{align}
where $S_t$, for example, is years in schooling for individual~$t$. Assume that
$\E(\varepsilon_{tj}) = 0$ and $\E(IQ_t \varepsilon_{tj}) = 0$ for $j = 1, 2$. Assume also that there is a
variable $MED$ (mother's education), which does not appear in either equation
but which is predetermined in that $\E(MED_t \varepsilon_{tj}) = 0$ for $j = 1, 2$. So the
common set of instruments is $\mathbf{x}_t = (1, IQ_t, MED_t)'$. The rest of the nonerror
variables in the system are endogenous variables. This two-equation system
has \emph{three} endogenous variables, $LW$, $S$, and $KWW$. Provided that the rank
condition is satisfied, the first equation is just identified because the number of
right-hand side variables equals the number of instruments, while the second
equation is overidentified.
\end{example}

\subsection*{The Complete System of Simultaneous Equations}

The ML counterpart of 3SLS is called the \textbf{full-information maximum likelihood
(FIML) estimator}. If it is to be estimated by FIML, the multiple-equation model
needs to be further specialized \emph{before} introducing the normality and i.i.d.\ assumptions. The required specialization is that those $M$ equations form a ``complete''
system of simultaneous equations. Completeness requires two things.

\begin{enumerate}
\item There are as many endogenous variables as there are equations. This condition
  implies that, if $\mathbf{y}_t$ is an $M$-dimensional vector collecting those $M$ endogenous
  variables, then $(y_{t1}, \ldots, y_{tM}, \mathbf{z}_{t1}, \ldots, \mathbf{z}_{tM})$ are all elements of
  $(\mathbf{y}_t, \mathbf{x}_t)$, which
  allows one to write the $M$-equation system \eqref{eq:8.5.1} as
  \begin{equation}
    \underset{(M \times M)}{\boldsymbol{\Gamma}_0}
    \underset{(M \times 1)}{\mathbf{y}_t}
    + \underset{(M \times K)}{\mathbf{B}_0}
    \underset{(K \times 1)}{\mathbf{x}_t}
    = \underset{(M \times 1)}{\boldsymbol{\varepsilon}_t}
    \quad (t = 1, 2, \ldots, n),
    \label{eq:8.5.8} \tag{8.5.8}
  \end{equation}
  with the $m$-th equation being
  $y_{tm} = \mathbf{z}_{tm}'\boldsymbol{\delta}_{0m} + \varepsilon_{tm}$. (This is illustrated for
  Example~8.1 below.) The equation system written in this way is called the
  \textbf{structural form}, and $(\boldsymbol{\Gamma}_0, \mathbf{B}_0)$ are called the \textbf{structural form parameters}. As
  will be illustrated below, each of the structural form parameters is a function of
  $(\boldsymbol{\delta}_{01}, \ldots, \boldsymbol{\delta}_{0M})$.

  The system in Example~8.1 satisfies this first requirement. The system in
  Example~8.2 is not a complete system because the number of endogenous variables is greater than the number of equations. It is not possible to estimate
  incomplete systems by FIML (unless we complete the system by adding appropriate equations, see the discussion about LIML below). This is in contrast to
  GMM; incomplete systems can be estimated by 3SLS (or more generally by
  multiple-equation GMM if conditional homoskedasticity is not assumed), as
  long as they satisfy the rank condition for identification.

\item The square matrix $\boldsymbol{\Gamma}_0$ is nonsingular. This implies that the structural form can
  be solved for the endogenous variables $\mathbf{y}_t$ as
  \begin{equation}
    \mathbf{y}_t = -\boldsymbol{\Gamma}_0^{-1}\mathbf{B}_0\mathbf{x}_t
    + \boldsymbol{\Gamma}_0^{-1}\boldsymbol{\varepsilon}_t
    \equiv \boldsymbol{\Pi}_0'\mathbf{x}_t + \mathbf{v}_t,
    \label{eq:8.5.9} \tag{8.5.9}
  \end{equation}
  where
  \begin{align}
    \underset{(M \times K)}{\boldsymbol{\Pi}_0'} &\equiv
    -\underset{(M \times M)}{\boldsymbol{\Gamma}_0^{-1}}
    \underset{(M \times K)}{\mathbf{B}_0},
    \label{eq:8.5.10} \tag{8.5.10} \\
    \underset{(M \times 1)}{\mathbf{v}_t} &\equiv
    \underset{(M \times M)}{\boldsymbol{\Gamma}_0^{-1}}
    \underset{(M \times 1)}{\boldsymbol{\varepsilon}_t}.
    \label{eq:8.5.11} \tag{8.5.11}
  \end{align}
  Expression \eqref{eq:8.5.9} is known as the \textbf{reduced-form} representation of the structural form, and the elements of $\boldsymbol{\Pi}_0$ are called the \textbf{reduced-form coefficients}. In
  each equation of the reduced form, the regressors are the same set of variables,
  $\mathbf{x}_t$. From \eqref{eq:8.5.3} and \eqref{eq:8.5.11}, it follows that
  $\E(\mathbf{x}_t\mathbf{v}_t') = \mathbf{0}$, namely, that all the
  regressors are predetermined. The reduced form, therefore, is a multivariate
  regression model.
\end{enumerate}

\subsection*{Relationship between \texorpdfstring{$(\boldsymbol{\Gamma}_0, \mathbf{B}_0)$}{(Gamma0, B0)} and \texorpdfstring{$\boldsymbol{\delta}_0$}{delta0}}

Let $\boldsymbol{\delta}_0$ be the stacked vector collecting all the coefficients in the $M$-equation system
\eqref{eq:8.5.1}:
\begin{equation}
  \underset{(\sum_{m=1}^{M} L_m \times 1)}{\boldsymbol{\delta}_0}
  = \begin{bmatrix}
      \boldsymbol{\delta}_{01} \\ \boldsymbol{\delta}_{02} \\ \vdots \\ \boldsymbol{\delta}_{0M}
    \end{bmatrix}.
  \label{eq:8.5.12} \tag{8.5.12}
\end{equation}
For understanding the mechanics of FIML, it is important to see how the coefficient matrices $(\boldsymbol{\Gamma}_0, \mathbf{B}_0)$ depend on $\boldsymbol{\delta}_0$. To illustrate, consider Example~8.1.
The $\boldsymbol{\delta}_0$ vector is
\begin{equation}
  \underset{(6 \times 1)}{\boldsymbol{\delta}_0}
  = \begin{bmatrix}
      \gamma_{11} \\ \beta_{11} \\ \beta_{12} \\
      \gamma_{21} \\ \beta_{21} \\ \beta_{22}
    \end{bmatrix}.
  \label{eq:8.5.13} \tag{8.5.13}
\end{equation}
The two endogenous variables are $(p_t, q_t)$. It does not matter how the endogenous
variables are ordered, so arbitrarily put $p_t$ first in $\mathbf{y}_t$:
$\mathbf{y}_t = (p_t, q_t)'$. The structural
form parameters can be written as
\begin{equation}
  \underset{(2 \times 2)}{\boldsymbol{\Gamma}_0}
  = \begin{bmatrix} -\gamma_{11} & 1 \\ -\gamma_{21} & 1 \end{bmatrix}, \quad
  \underset{(2 \times 3)}{\mathbf{B}_0}
  = \begin{bmatrix}
      -\beta_{11} & -\beta_{12} & 0 \\
      -\beta_{21} & 0 & -\beta_{22}
    \end{bmatrix}.
  \label{eq:8.5.14} \tag{8.5.14}
\end{equation}

This example serves to illustrate three points. First, each row of $\boldsymbol{\Gamma}_0$ has an
element of unity, reflecting the fact that the coefficient of the dependent variable
in each of the $M$ equations \eqref{eq:8.5.1} is unity. In this sense $\boldsymbol{\Gamma}_0$ is already normalized. Second, some of the elements of $\boldsymbol{\Gamma}_0$ and $\mathbf{B}_0$ are zeros, reflecting the fact that some
endogenous or predetermined variables do not appear in certain equations of the
$M$-equation system. This feature of the structural form coefficient matrices is
called the \textbf{exclusion restrictions}. (In Example~8.1, no elements of $\boldsymbol{\Gamma}_0$ happen
to be zero because the two endogenous variables appear in both equations.) Third,
the system involves no cross-equation restrictions, so each element of $\boldsymbol{\delta}_0$ appears
in $(\boldsymbol{\Gamma}_0, \mathbf{B}_0)$ just once.

\subsection*{The FIML Likelihood Function}

To make the complete system of simultaneous equations estimable by FIML, we
assume (1) that the structural error vector $\boldsymbol{\varepsilon}_t$ is jointly normal conditional on $\mathbf{x}_t$,
i.e., $\boldsymbol{\varepsilon}_t \mid \mathbf{x}_t \sim N(\mathbf{0}, \boldsymbol{\Sigma}_0)$, and
(2) that $\{\mathbf{y}_t, \mathbf{x}_t\}$ is i.i.d., not just ergodic stationary
martingale differences (this assumption strengthens Assumptions~4.2 and~4.5).

By the normality assumption~(1) about the structural errors $\boldsymbol{\varepsilon}_t$ and
\eqref{eq:8.5.11},
we have $\mathbf{v}_t \mid \mathbf{x}_t \sim N(\mathbf{0},
\boldsymbol{\Gamma}_0^{-1}\boldsymbol{\Sigma}_0(\boldsymbol{\Gamma}_0^{-1})')$.
Combining this and the reduced form
$\mathbf{y}_t = \boldsymbol{\Pi}_0'\mathbf{x}_t + \mathbf{v}_t$ with \eqref{eq:8.5.10} produces
\begin{equation}
  \mathbf{y}_t \mid \mathbf{x}_t \sim
  N\bigl(-\boldsymbol{\Gamma}_0^{-1}\mathbf{B}_0\mathbf{x}_t,\;
    \boldsymbol{\Gamma}_0^{-1}\boldsymbol{\Sigma}_0(\boldsymbol{\Gamma}_0^{-1})'\bigr).
  \label{eq:8.5.15} \tag{8.5.15}
\end{equation}
So the log conditional likelihood for observation~$t$ is
\begin{multline}
  \log f(\mathbf{y}_t \mid \mathbf{x}_t;\, \boldsymbol{\delta}, \boldsymbol{\Sigma})
  = -\frac{M}{2}\log(2\pi)
    - \frac{1}{2}\log(|\boldsymbol{\Gamma}^{-1}\boldsymbol{\Sigma}(\boldsymbol{\Gamma}^{-1})'|)
    \\
    - \frac{1}{2}[\mathbf{y}_t + \boldsymbol{\Gamma}^{-1}\mathbf{B}\mathbf{x}_t]'
      [\boldsymbol{\Gamma}^{-1}\boldsymbol{\Sigma}(\boldsymbol{\Gamma}^{-1})']^{-1}
      [\mathbf{y}_t + \boldsymbol{\Gamma}^{-1}\mathbf{B}\mathbf{x}_t].
  \label{eq:8.5.16} \tag{8.5.16}
\end{multline}
The likelihood is a function of $(\boldsymbol{\delta}, \boldsymbol{\Sigma})$ because, as illustrated in \eqref{eq:8.5.14}, the structural form coefficients $(\boldsymbol{\Gamma}, \mathbf{B})$ are functions of $\boldsymbol{\delta}$. Now
\begin{equation}
  [\mathbf{y}_t + \boldsymbol{\Gamma}^{-1}\mathbf{B}\mathbf{x}_t]'
  [\boldsymbol{\Gamma}^{-1}\boldsymbol{\Sigma}(\boldsymbol{\Gamma}^{-1})']^{-1}
  [\mathbf{y}_t + \boldsymbol{\Gamma}^{-1}\mathbf{B}\mathbf{x}_t]
  = [\boldsymbol{\Gamma}\mathbf{y}_t + \mathbf{B}\mathbf{x}_t]'
    \boldsymbol{\Sigma}^{-1}
    [\boldsymbol{\Gamma}\mathbf{y}_t + \mathbf{B}\mathbf{x}_t]
  \label{eq:8.5.17} \tag{8.5.17}
\end{equation}
and
\begin{equation}
  |\boldsymbol{\Gamma}^{-1}\boldsymbol{\Sigma}(\boldsymbol{\Gamma}^{-1})'|
  = |\boldsymbol{\Gamma}^{-1}| \,|\boldsymbol{\Sigma}|\, |(\boldsymbol{\Gamma}^{-1})'|
  = |\boldsymbol{\Sigma}|/|\boldsymbol{\Gamma}|^2.
  \label{eq:8.5.18} \tag{8.5.18}
\end{equation}
Substituting these into \eqref{eq:8.5.16},
\begin{multline}
  \log f(\mathbf{y}_t \mid \mathbf{x}_t;\, \boldsymbol{\delta}, \boldsymbol{\Sigma})
  = -\frac{M}{2}\log(2\pi)
    + \frac{1}{2}\log(|\boldsymbol{\Gamma}|^2)
    - \frac{1}{2}\log(|\boldsymbol{\Sigma}|) \\
    - \frac{1}{2}[\boldsymbol{\Gamma}\mathbf{y}_t + \mathbf{B}\mathbf{x}_t]'
      \boldsymbol{\Sigma}^{-1}
      [\boldsymbol{\Gamma}\mathbf{y}_t + \mathbf{B}\mathbf{x}_t].
  \label{eq:8.5.19} \tag{8.5.19}
\end{multline}
Taking the average over~$t$, we obtain the FIML objective function for a random
sample:
\begin{multline}
  Q_n(\boldsymbol{\delta}, \boldsymbol{\Sigma})
  = -\frac{M}{2}\log(2\pi)
    + \frac{1}{2}\log(|\boldsymbol{\Gamma}|^2)
    - \frac{1}{2}\log(|\boldsymbol{\Sigma}|) \\
    - \frac{1}{2n}\sum_{t=1}^{n}
      [\boldsymbol{\Gamma}\mathbf{y}_t + \mathbf{B}\mathbf{x}_t]'
      \boldsymbol{\Sigma}^{-1}
      [\boldsymbol{\Gamma}\mathbf{y}_t + \mathbf{B}\mathbf{x}_t].
  \label{eq:8.5.20} \tag{8.5.20}
\end{multline}
The FIML estimate of $(\boldsymbol{\delta}_0, \boldsymbol{\Sigma}_0)$ is the $(\boldsymbol{\delta}, \boldsymbol{\Sigma})$ that maximizes this objective function.

\subsection*{The FIML Concentrated Likelihood Function}

As in the maximization of the objective function of the multivariate regression
model, we can proceed in two steps. The first step is to maximize
$Q_n(\boldsymbol{\delta}, \boldsymbol{\Sigma})$ with
respect to $\boldsymbol{\Sigma}$ given $\boldsymbol{\delta}$. This yields
\begin{equation}
  \underset{(M \times M)}{\hat{\boldsymbol{\Sigma}}(\boldsymbol{\delta})}
  = \frac{1}{n}\sum_{t=1}^{n}
    (\boldsymbol{\Gamma}\mathbf{y}_t + \mathbf{B}\mathbf{x}_t)
    (\boldsymbol{\Gamma}\mathbf{y}_t + \mathbf{B}\mathbf{x}_t)'.
  \label{eq:8.5.21} \tag{8.5.21}
\end{equation}
Its $(m, h)$ element can be written as
\begin{equation}
  \hat{\sigma}_{mh}
  = \frac{1}{n}\sum_{t=1}^{n}
    (y_{tm} - \mathbf{z}_{tm}'\boldsymbol{\delta}_m)
    (y_{th} - \mathbf{z}_{th}'\boldsymbol{\delta}_h).
  \label{eq:8.5.22} \tag{8.5.22}
\end{equation}
By substituting \eqref{eq:8.5.21} back into the objective function \eqref{eq:8.5.20},
$(1/n$ times) the
FIML concentrated likelihood function can be derived as
\begin{align}
  Q_n^*(\boldsymbol{\delta})
  &\equiv Q_n(\boldsymbol{\delta},\, \hat{\boldsymbol{\Sigma}}(\boldsymbol{\delta})) \notag \\
  &= -\frac{M}{2}\log(2\pi)
     - \frac{M}{2}
     + \frac{1}{2}\log(|\boldsymbol{\Gamma}|^2)
     - \frac{1}{2}\log(|\hat{\boldsymbol{\Sigma}}(\boldsymbol{\delta})|) \notag \\
  &= -\frac{M}{2}\log(2\pi)
     - \frac{M}{2}
     - \frac{1}{2}\log\biggl|\frac{1}{n}\sum_{t=1}^{n}
       (\mathbf{y}_t + \boldsymbol{\Gamma}^{-1}\mathbf{B}\mathbf{x}_t)
       (\mathbf{y}_t + \boldsymbol{\Gamma}^{-1}\mathbf{B}\mathbf{x}_t)'\biggr|.
  \label{eq:8.5.23} \tag{8.5.23}
\end{align}

The second step is to maximize this FIML concentrated likelihood with respect to
$\boldsymbol{\delta}$. The FIML estimator $\hat{\boldsymbol{\delta}}$ of $\boldsymbol{\delta}_0$ is the $\boldsymbol{\delta}$ that maximizes this function.\footnote{So the FIML estimator $\hat{\boldsymbol{\delta}}$ is an extremum estimator that maximizes
$-\bigl|\frac{1}{n}\sum_{t=1}^{n}
(\mathbf{y}_t + \boldsymbol{\Gamma}^{-1}\mathbf{B}\mathbf{x}_t)
(\mathbf{y}_t + \boldsymbol{\Gamma}^{-1}\mathbf{B}\mathbf{x}_t)'\bigr|$.
Some aspects of this extremum estimator are explored in an optional analytical exercise (see Analytical Exercise~3).}
The FIML
estimator of $\boldsymbol{\Sigma}_0$ is $\hat{\boldsymbol{\Sigma}}(\hat{\boldsymbol{\delta}})$.

\subsection*{Testing Overidentifying Restrictions}

Now compare this FIML concentrated likelihood $Q_n^*(\boldsymbol{\delta})$ given in \eqref{eq:8.5.23} with that
for the multivariate regression model, $Q_n^*(\boldsymbol{\Pi})$ given in \eqref{eq:8.4.13}. The FIML concentrated likelihood $Q_n^*(\boldsymbol{\delta})$ is what obtains when we impose on
$Q_n^*(\boldsymbol{\Pi})$ the restrictions
implied by the null hypothesis that
\begin{equation}
  \mathrm{H}_0:\;
  \boldsymbol{\Pi}_0' = -\boldsymbol{\Gamma}_0^{-1}\mathbf{B}_0
  \;\text{ or }\;
  \boldsymbol{\Gamma}_0\boldsymbol{\Pi}_0' + \mathbf{B}_0 = 0.
  \label{eq:8.5.24} \tag{8.5.24}
\end{equation}
Put differently, FIML estimation of $\boldsymbol{\delta}_0$ is a constrained multivariate regression.
Therefore, the truth of the null hypothesis can be tested by the likelihood ratio
principle. Since the OLS estimator $\hat{\boldsymbol{\Pi}}$ given in \eqref{eq:8.4.5} maximizes the multivariate
regression concentrated likelihood function $Q_n^*(\boldsymbol{\Pi})$ given in \eqref{eq:8.4.13} and the FIML
estimator $\hat{\boldsymbol{\delta}}$ maximizes the FIML concentrated likelihood function $Q_n^*(\boldsymbol{\delta})$ given in
\eqref{eq:8.5.23}, the likelihood ratio test statistic is
\begin{align}
  LR &= 2n \cdot (Q_n^*(\hat{\boldsymbol{\Pi}}) - Q_n^*(\hat{\boldsymbol{\delta}})) \notag \\
     &= n \times \bigl(\log|\hat{\boldsymbol{\Omega}}(-(\tilde{\boldsymbol{\Gamma}}^{-1}\tilde{\mathbf{B}})' )|
        - \log|\hat{\boldsymbol{\Omega}}(\hat{\boldsymbol{\Pi}})|\bigr),
  \label{eq:8.5.25} \tag{8.5.25}
\end{align}
where $\hat{\boldsymbol{\Omega}}(\boldsymbol{\Pi})$ is defined in \eqref{eq:8.4.10} and
$(\tilde{\boldsymbol{\Gamma}}, \tilde{\mathbf{B}})$ is the value of $(\boldsymbol{\Gamma}, \mathbf{B})$ implied by $\hat{\boldsymbol{\delta}}$.
Since the dimension of $\boldsymbol{\delta}$ equals $\sum_{m=1}^{M} L_m$, the number of restrictions implied by
the null hypothesis $\mathrm{H}_0$ above is
\begin{equation}
  KM - \sum_{m=1}^{M} L_m = \sum_{m=1}^{M}(K - L_m),
  \label{eq:8.5.26} \tag{8.5.26}
\end{equation}
which is the total number of overidentifying restrictions of the system. Therefore,
the likelihood ratio test based on the $LR$ statistic \eqref{eq:8.5.25} is a \textbf{test of overidentifying restrictions}. It is a specification test because the restriction being tested is a
condition assumed in the model.

\subsection*{Properties of the FIML Estimator}

There is no closed-form solution to the FIML estimator. So, in order to establish
the consistency and asymptotic normality, we would have to verify the conditions
of relevant theorems of the previous chapter. We summarize here, without proof,
the main properties of the FIML estimator.

\subsubsection*{Identification}

Stating the identification condition in various forms for complete simultaneous
equations systems is a major topic in most textbooks. Here it is left as an optional
analytical exercise (Analytical Exercise~3) to show that the identification condition for FIML as an extremum estimator is equivalent to the rank condition for
identification (that $\E(\mathbf{x}_t\mathbf{z}_{tm}')$ be of full column rank for all
$m = 1, 2, \ldots, M$).

\subsubsection*{Asymptotic Properties}

Although the FIML objective function \eqref{eq:8.5.20} is derived under the normality of
$\boldsymbol{\varepsilon}_t \mid \mathbf{x}_t$, the FIML estimator of $\boldsymbol{\delta}_0$ is consistent and asymptotically normal without
the normality assumption. For a proof of consistency without normality, see, e.g.,
Amemiya (1985, pp.~232--233). There is a nice proof of asymptotic normality
of $\hat{\boldsymbol{\delta}}$ due to Hausman (1975), which shows that an iterative instrumental variable
estimation can be viewed as an algorithm for solving the first-order conditions for
the maximization of $Q_n(\boldsymbol{\delta}, \boldsymbol{\Sigma})$.\footnote{We will not display the first-order conditions (the likelihood equations) for FIML because it requires
cumbersome matrix notation. See Amemiya (1985, pp.~233--234) or Davidson and MacKinnon (1993, pp.~658--660) for more details. For the SUR model, the likelihood equations and the iterative procedure are much easier
to describe; see the next subsection.}
This argument shows that the FIML estimator of
$\boldsymbol{\delta}_0$ is asymptotically equivalent to the 3SLS estimator.\footnote{This asymptotic equivalence does not carry over to nonlinear equation systems; nonlinear FIML is more
efficient than nonlinear 3SLS. See Amemiya (1985, Section~8.2).}
Therefore, its asymptotic
variance is given by (4.5.15), and a consistent estimator of the asymptotic variance
is (4.5.17).

\subsubsection*{Invariance}

Since the FIML estimator is an ML estimator, it has the invariance property discussed in Section~7.1. To illustrate, consider Example~8.1 and suppose you wish to
reparameterize the supply equation as
\begin{equation}
  p_t = \tilde{\gamma}_{21} q_t + \tilde{\beta}_{21}
  + \tilde{\beta}_{22} w_t + \tilde{\varepsilon}_{t2}
  \quad \text{(supply)}.
  \label{eq:8.5.27} \tag{8.5.27}
\end{equation}
The old supply parameters $(\gamma_{21}, \beta_{21}, \beta_{22})$ and the new supply parameters
$(\tilde{\gamma}_{21}, \tilde{\beta}_{21}, \tilde{\beta}_{22})$ are connected by the one-to-one mapping
\begin{equation}
  \tilde{\gamma}_{21} = 1/\gamma_{21}, \quad
  \tilde{\beta}_{21} = -\beta_{21}/\gamma_{21}, \quad
  \tilde{\beta}_{22} = -\beta_{22}/\gamma_{21}.
  \label{eq:8.5.28} \tag{8.5.28}
\end{equation}
If $(\hat{\gamma}_{21}, \hat{\beta}_{21}, \hat{\beta}_{22})$ is the FIML estimator of
$(\gamma_{21}, \beta_{21}, \beta_{22})$ for the system consisting
of the demand equation \eqref{eq:8.5.4} and the old supply equation \eqref{eq:8.5.5}, then the FIML
estimator based on the system consisting of \eqref{eq:8.5.4} and the new supply equation
\eqref{eq:8.5.27}
is numerically the same as the value of the above mapping at
$(\hat{\gamma}_{21}, \hat{\beta}_{21}, \hat{\beta}_{22})$. The
3SLS estimator (and 2SLS for that matter) does not have this invariance property.

The foregoing discussion about FIML can be summarized as

\begin{proposition}[Asymptotic properties of FIML]\label{prop:8.1}
Consider the $M$-equation system \eqref{eq:8.5.1} and let $\boldsymbol{\delta}_0$ be the stacked vector collecting all the coefficients in the
$M$-equation system. Make the following assumptions.
\begin{itemize}
\item the rank condition for identification (that $\E(\mathbf{x}_t\mathbf{z}_{tm}')$ be of full column rank for all
  $m = 1, 2, \ldots, M$) is satisfied,
\item $\E(\mathbf{x}_t\mathbf{x}_t')$ is nonsingular,
\item the $M$-equation system can be written as a complete system of simultaneous
  equations \eqref{eq:8.5.8} with $\boldsymbol{\Gamma}_0$ nonsingular,
\item $\boldsymbol{\varepsilon}_t \mid \mathbf{x}_t \sim N(\mathbf{0}, \boldsymbol{\Sigma}_0)$,
  $\boldsymbol{\Sigma}_0$ positive definite,
\item $\{\mathbf{y}_t, \mathbf{x}_t\}$ is i.i.d.,
\item the parameter space for $(\boldsymbol{\delta}_0, \boldsymbol{\Sigma}_0)$ is compact with the true parameter vector
  $(\boldsymbol{\delta}_0, \boldsymbol{\Sigma}_0)$ included as an interior point.
\end{itemize}
Then
\begin{enumerate}[label=(\alph*)]
\item the FIML estimator $(\hat{\boldsymbol{\delta}}, \hat{\boldsymbol{\Sigma}})$, which maximizes \eqref{eq:8.5.20}, is consistent and asymptotically normal,
\item the asymptotic variance of $\hat{\boldsymbol{\delta}}$ is given by $(4.5.15)$, and a consistent estimator of
  the asymptotic variance is $(4.5.17)$ where $\hat{\sigma}^{mh}$ is the $(m, h)$ element of
  $\hat{\boldsymbol{\Sigma}}^{-1}$,
\item the likelihood ratio statistic \eqref{eq:8.5.25} for testing overidentifying restrictions is
  asymptotically $\chi^2$ with $KM - \sum_{m=1}^{M} L_m$ degrees of freedom.
\end{enumerate}
Furthermore, these asymptotic results about $\hat{\boldsymbol{\delta}}$ hold even if
$\boldsymbol{\varepsilon}_t \mid \mathbf{x}_t$ is not normal.
\end{proposition}

\subsection*{ML Estimation of the SUR Model}

The SUR model is a specialization of the FIML model with $\mathbf{z}_{tm}$ being a subvector
of $\mathbf{x}_t$. It is therefore an $M$-equation system \eqref{eq:8.5.1} where the regressors $\mathbf{z}_{tm}$ are
predetermined. Unlike in the multivariate regression model, the set of regressors
could differ across equations. It is easy to show (see Review Question~3) that no
two dependent variables can be the same variable, so $\mathbf{y}_t$ ($M \times 1$) in the structural
form \eqref{eq:8.5.8} simply collects the $M$ dependent variables. Furthermore, since $\mathbf{z}_{tm}$
includes no endogenous variables, $\boldsymbol{\Gamma}_0$ in the structural form is the identity matrix.
Thus, the FIML objective function \eqref{eq:8.5.20} becomes
\begin{equation}
  Q_n(\boldsymbol{\delta}, \boldsymbol{\Sigma})
  = -\frac{M}{2}\log(2\pi)
    - \frac{1}{2}\log(|\boldsymbol{\Sigma}|)
    - \frac{1}{2n}\sum_{t=1}^{n}
      [\boldsymbol{\Gamma}\mathbf{y}_t + \mathbf{B}\mathbf{x}_t]'
      \boldsymbol{\Sigma}^{-1}
      [\boldsymbol{\Gamma}\mathbf{y}_t + \mathbf{B}\mathbf{x}_t],
  \label{eq:8.5.29} \tag{8.5.29}
\end{equation}
where $\boldsymbol{\Gamma}$ is understood to be the identity matrix. (We continue to carry $\boldsymbol{\Gamma}$ around
in order to facilitate the comparison to the objective function for LIML to be introduced in the next subsection.) As in FIML, the $\boldsymbol{\Sigma}$ that maximizes the objective
function given $\boldsymbol{\delta}$ is given by $\hat{\boldsymbol{\Sigma}}(\boldsymbol{\delta})$ in \eqref{eq:8.5.21}.

For the SUR model, the score with respect to $\boldsymbol{\delta}$ can be written down rather
easily. The $m$-th row of $\boldsymbol{\Gamma}\mathbf{y}_t + \mathbf{B}\mathbf{x}_t$ is
$y_{tm} - \mathbf{z}_{tm}'\boldsymbol{\delta}_m$. Define
$\mathbf{y}$ ($Mn \times 1$) and $\mathbf{Z}$
($Mn \times \sum_{m=1}^{M} L_m$) as in the first analytical exercise to Chapter~4. Then
\begin{equation}
  \sum_{t=1}^{n}
    [\boldsymbol{\Gamma}\mathbf{y}_t + \mathbf{B}\mathbf{x}_t]'
    \boldsymbol{\Sigma}^{-1}
    [\boldsymbol{\Gamma}\mathbf{y}_t + \mathbf{B}\mathbf{x}_t]
  = (\mathbf{y} - \mathbf{Z}\boldsymbol{\delta})'
    (\boldsymbol{\Sigma}^{-1} \otimes \mathbf{I}_n)
    (\mathbf{y} - \mathbf{Z}\boldsymbol{\delta}).
  \label{eq:8.5.30} \tag{8.5.30}
\end{equation}
Therefore,
\begin{equation}
  \frac{\partial Q_n(\boldsymbol{\delta}, \boldsymbol{\Sigma})}{\partial \boldsymbol{\delta}}
  = \frac{1}{n}\,\mathbf{Z}'(\boldsymbol{\Sigma}^{-1} \otimes \mathbf{I}_n)
    (\mathbf{y} - \mathbf{Z}\boldsymbol{\delta}).
  \label{eq:8.5.31} \tag{8.5.31}
\end{equation}
(Doing the same for the more general case of FIML is not as easy because, although
\eqref{eq:8.5.30} holds for FIML as well, the score $\frac{\partial Q_n}{\partial \boldsymbol{\delta}}$ is more complicated thanks to the
term $\log(|\boldsymbol{\Gamma}|^2)$ in the FIML likelihood function. Put differently, if $|\boldsymbol{\Gamma}|$ is a constant
(i.e., independent of $\boldsymbol{\delta}$), then $\frac{\partial Q_n}{\partial \boldsymbol{\delta}}$ for FIML is the same as \eqref{eq:8.5.31}.)
Setting this
equal to zero and solving for $\boldsymbol{\delta}$, we obtain
\begin{equation}
  \hat{\boldsymbol{\delta}}(\boldsymbol{\Sigma})
  \equiv [\mathbf{Z}'(\boldsymbol{\Sigma}^{-1} \otimes \mathbf{I}_n)\mathbf{Z}]^{-1}
    \mathbf{Z}'(\boldsymbol{\Sigma}^{-1} \otimes \mathbf{I}_n)\mathbf{y}.
  \label{eq:8.5.32} \tag{8.5.32}
\end{equation}
Given some consistent estimator $\hat{\boldsymbol{\Sigma}}$ of $\boldsymbol{\Sigma}_0$,
$\hat{\boldsymbol{\delta}}(\hat{\boldsymbol{\Sigma}})$ is the SUR estimator of $\boldsymbol{\delta}_0$ derived
in Chapter~4.

These two functions $(\hat{\boldsymbol{\delta}}(\boldsymbol{\Sigma}),\, \hat{\boldsymbol{\Sigma}}(\boldsymbol{\delta}))$ define a mapping from the parameter space
for $(\boldsymbol{\delta}, \boldsymbol{\Sigma})$ to itself, and a solution to the first-order conditions for the maximization
of $Q_n(\boldsymbol{\delta}, \boldsymbol{\Sigma})$ is a fixed point in this mapping. Such a fixed point can be calculated
by the following \textbf{iterative SUR}. Let
$(\hat{\boldsymbol{\delta}}^{(j)}, \hat{\boldsymbol{\Sigma}}^{(j)})$ be the estimate of
$(\boldsymbol{\delta}_0, \boldsymbol{\Sigma}_0)$ in the
$j$-th iteration. The estimate in the next iteration is
\begin{equation}
  \hat{\boldsymbol{\delta}}^{(j+1)} = \hat{\boldsymbol{\delta}}(\hat{\boldsymbol{\Sigma}}^{(j)}), \quad
  \hat{\boldsymbol{\Sigma}}^{(j+1)} = \hat{\boldsymbol{\Sigma}}(\hat{\boldsymbol{\delta}}^{(j+1)}).
  \label{eq:8.5.33} \tag{8.5.33}
\end{equation}
The first part of this iteration is the SUR estimation given $\hat{\boldsymbol{\Sigma}}^{(j)}$, while the latter
part updates the estimate of $\boldsymbol{\Sigma}_0$ using the current estimate of $\boldsymbol{\delta}_0$. If this process
converges, then the limit is a solution to the first-order conditions.

\bigskip
\noindent\rule{\textwidth}{0.4pt}
\textbf{QUESTIONS FOR REVIEW}
\medskip

\begin{enumerate}
\item (Deriving $f(\mathbf{y}_t \mid \mathbf{x}_t)$ from $f(\boldsymbol{\varepsilon}_t \mid \mathbf{x}_t)$)\quad
  In this exercise, we wish to derive the
  likelihood for observation~$t$ for the FIML model by the following fact from
  probability theory:

  \begin{quote}
  \textit{Change of Variable Formula:}\quad Suppose $\boldsymbol{\varepsilon}$ is a continuous $M$-dimensional random vector with density $f_{\boldsymbol{\varepsilon}}(\boldsymbol{\varepsilon})$ and
  $\mathbf{g}: \mathbb{R}^M \to \mathbb{R}^M$ is a one-to-one and differentiable mapping on an open set~$S$, with the inverse
  mapping denoted as $\mathbf{g}^{-1}(\cdot)$. Let the $M \times M$ matrix $\mathbf{J}(\boldsymbol{\varepsilon})$ be the Jacobian
  matrix $\frac{\partial \mathbf{g}(\boldsymbol{\varepsilon})}{\partial \boldsymbol{\varepsilon}'}$
  (the $M \times M$ matrix of partial derivatives, whose $(i, j)$
  element is $\frac{\partial g_i(\boldsymbol{\varepsilon})}{\partial \varepsilon_j}$).
  Assume that $|\mathbf{J}(\boldsymbol{\varepsilon})|$ is not zero at any point in~$S$.
  Then the density of $\mathbf{y} = \mathbf{g}(\boldsymbol{\varepsilon})$ is given by
  \[
    f(\mathbf{y}) = \frac{f_{\boldsymbol{\varepsilon}}(\mathbf{g}^{-1}(\mathbf{y}))}
    {\text{abs}(|\mathbf{J}(\mathbf{g}^{-1}(\mathbf{y}))|)}.
  \]
  \end{quote}

  (Note that $|\mathbf{J}(\mathbf{g}^{-1}(\mathbf{y}))|$ is a determinant, which may or may not be positive.)
  Using this fact, derive the density of $\mathbf{y}_t \mid \mathbf{x}_t$ from the density of
  $\boldsymbol{\varepsilon}_t \mid \mathbf{x}_t$
  and verify that the associated log likelihood is given by \eqref{eq:8.5.16}.
  \textbf{Hint:} The
  inverse mapping $\mathbf{g}^{-1}(\mathbf{y}_t)$ is the left-hand side of \eqref{eq:8.5.8}. So
  $\mathbf{J}(\boldsymbol{\varepsilon}_t) = \boldsymbol{\Gamma}_0^{-1}$. Also,
  $\log(\text{abs}(|\boldsymbol{\Gamma}|)) = \frac{1}{2}\log(|\boldsymbol{\Gamma}|^2)$.

\item (Instruments outside the system)\quad
  Consider the complete system of simultaneous equations \eqref{eq:8.5.8} with $|\boldsymbol{\Gamma}_0| \neq 0$. Suppose that one of the $K$ predetermined
  variables, say $x_{tK}$, does not appear in any of the $M$-equations. Let
  $\hat{\mathbf{E}}^*(\mathbf{y} \mid \mathbf{x})$ be
  the least squares projection of $\mathbf{y}$ on $\mathbf{x}$. Show that
  $\hat{\mathbf{E}}^*(y_{tm} \mid \mathbf{x}_t) = \hat{\mathbf{E}}^*(y_{tm} \mid \tilde{\mathbf{x}}_t)$,
  where $\tilde{\mathbf{x}}_t = (x_{t1}, \ldots, x_{t,K-1})'$ (so
  $\mathbf{x}_t = (\tilde{\mathbf{x}}_t', x_{tK})'$).
  (As you will show in Analytical Exercise~4, including $x_{tK}$ in addition to $\tilde{\mathbf{x}}_t$ in the list of instruments will
  not change the asymptotic variance of the FIML estimator of $\boldsymbol{\delta}_0$; it probably
  makes the small sample properties only worse.)

\item Recall that in Example~8.1 $y_{t1}$ and $y_{t2}$ are the same variable. In the SUR model,
  this cannot happen; that is, no two elements of $(y_{t1}, y_{t2}, \ldots, y_{tM})$ are the same
  variable if the assumptions describing the model are satisfied. Prove this.
  \textbf{Hint:}
  Suppose $y_{t1}$ and $y_{t2}$ are the same variable. Then
  \[
    \varepsilon_{t1} - \varepsilon_{t2}
    = \mathbf{z}_{t1}'\boldsymbol{\delta}_{0,1} - \mathbf{z}_{t2}'\boldsymbol{\delta}_{0,2}.
  \]
\end{enumerate}
